\usepackage[a4paper]{geometry} % 页面尺寸设置包(窄)
% 窄边：left=1.27cm,right=1.27cm,top=1.27cm,bottom=1.27cm
% 正常：left=2.5cm,right=2.5cm,top=2.54cm,bottom=2.54cm


\usepackage{amssymb}
\usepackage{amsfonts}
\usepackage{amsthm}
\usepackage{amsmath}
\usepackage{unicode-math} % 定义跳跃符号⟦需要(但是与bm包不兼容!)


\usepackage{extarrows} % 数学箭头包

\usepackage{booktabs} % 三线表

\usepackage{multirow} % 合并单元格


% \usepackage{bm} % 加粗字体包(与unicode-math不兼容)，使用\mathbf替代\bm更好
\def\bm#1{\symbf{#1}} % 当使用unicode-math时，这样加粗
% \def\bm#1{\mathbf{#1}} % 粗体(对希腊字母无效)


% \usepackage{cases} % 方程组每个方程编号，数学环境只能内嵌在它里面


\usepackage{graphicx} % 使用resizebox缩放


\usepackage{subcaption} % 并排表格、subfigure子图


\usepackage{esint} % 环路积分、通量符号


\usepackage{float} % 图片figure加入大写的H，可以强制位置


% \usepackage{enumerate} % 改变enumerate编号格式


\usepackage{mathrsfs} % 花体\mathscr，需要安装jknapltx(TexLive)或者jknappen(MikTex)包


\usepackage{hyperref} % 可以在前面加入[hidelinks]来不显示超链接框，提供各种超链接；在标题需要写数学公式时，用\texorpdfstring{Tex}{String}替代，其中Tex写数学公式，String写用于替代公式的文本串(用于在书签中显式)


% \usepackage{cite} % 引用参考文献，superscript使得引用编号是上标的形式


% \usepackage[title]{appendix} % 附录标题


\usepackage{enumitem} % 调整enumerate和enumerate的行间距以及编号格式


% \usepackage{algorithm} % 与 algorithm2e 不兼容
% \usepackage{algorithmic} % 与 algorithm2e 不兼容
\usepackage[ruled,linesnumbered]{algorithm2e}


% \usepackage{rotating}%用于旋转对象


\usepackage{longtable} % 跨页表格(请勿在外层使用center或者table)
% \usepackage{multirow} % 单元格纵向合并


%\usepackage{xcolor}%字体颜色


\usepackage[backend=bibtex, sorting=none]{biblatex} % 选择后端处理方式为bibtex(需要运行bibtex)，编号方式是数值编号，禁用重排序，引用格式是gb7714-2015(需要xstring.sty)，需要安装biblatex包(以及biblatex-gb7714-2015包)


\newcommand{\tabincell}[2]{\begin{tabular}{@{}#1@{}}#2\end{tabular}} % 表格单元格内换行
\newtheoremstyle{Problem}   % name
    {15pt}                  % Space above
    {5pt}                   % Space below
    {}                      % Body font
    {0em}                   % Indent amount
    {\bfseries}             % Theorem head font
    {:}                     % Punctuation after theorem head
    {0.5em}                 % Space after theorem head
    {}                      % Theorem head spec (can be left empty, meaning `normal' )
% 自定义环境样式，请查询amsthm
\newtheoremstyle{Solution}  % name
    {5pt}                   % Space above
    {7pt}                   % Space below
    {}                      % Body font
    {0em}                   % Indent amount
    {\ttfamily}             % Theorem head font
    {:}                     % Punctuation after theorem head
    {0.5em}                 % Space after theorem head
    {}                      % Theorem head spec (can be left empty, meaning `normal' )
% 自定义环境样式，请查询amsthm

\theoremstyle{Problem} % 以下自定义的环境使用已定义格式(正体(宋体)加粗)
    \newtheorem{assumption}{假设}[section]
    \newtheorem{proposition}{命题}[section]
    \newtheorem{remark}{注}[section]
    \newtheorem{algo}{算法}[section]
    \newtheorem{theorem}{定理}[section]
    \newtheorem{corollary}{推论}[section]
    \newtheorem{lemma}{引理}[section]
    \newtheorem{definition}{定义}[section]
    \newtheorem{example}{例子}[section]
    \newtheorem{axiom}{公理}[section]

% \theoremstyle{Solution}
%     \newtheorem*{solution}{解} % 解不加编号
%     \newtheorem*{Proof}{证明}

\def\R{\mathbb{R}}
\def\F{\mathbb{F}}
\def\C{\mathbb{C}}
\def\Q{\mathbb{Q}}
\def\d{{\rm d}}
\def\D{{\mathcal{D}}} % 微分算子
\def\Z{\mathbb{Z}}
\def\N{\mathbb{N}}
\def\X{\chi} % 特征函数，示性函数
\def\e{{\rm e}}
\def\i{{\rm i}}
\def\rank{\mathop{\rm rank}}
\def\tr{\mathop{\rm tr}}
\def\diam{{\rm diam}}
\def\diag{\mathop{\rm diag}}
\def\lcm{\mathop{\rm lcm}}
\def\mod{{\;\rm mod\,}}
\def\norm#1{\left\lVert#1\right\rVert} % 范数
\def\B#1{\left\{#1\right\}} % 大括号
\def\b#1{\left[#1\right]} % 中括号
\def\abs#1{\left\lvert#1\right\rvert}
\def\bdots{\begin{rotate}{90}$\ddots$\end{rotate}} % 反对角的省略号
\def\jump#1{\left\lBrack #1 \right\rBrack} % 跳跃
\def\avg#1{\left\lbrace #1 \right\rbrace} % 均值

\setmonofont{Consolas} % 设置等宽字体样式
\allowdisplaybreaks % 允许公式换页